% chapters/ch02h_passive_pattern_paper.tex
% Embedded Passive Pattern peer-review paper

\section{Passive Pattern --- Mechanistic Theory and Constrained
Stochastic Control of Dialogue (Embedded Peer-Review Paper)}
\label{sec:passive-pattern-paper}

\nrmobox{Paper status}{
\textbf{Title}: Passive Pattern --- Mechanistic Theory and
Constrained Stochastic Control of Dialogue\\[0.3em]
\textbf{Revision}: This revision adds a full ``why-it-works''
theory: variance amplification under partial observability,
feasible-action-set shrinkage, and irreversibility-driven tail risk.
}

\subsection{Contents}
\label{sec:pp-contents}

Abstract; 1.~Introduction; 2.~Claims and Scope; 3.~Related Work;
4.~Setup --- Dialogue as Non-Ergodic Controlled Process;
5.~Mechanistic Theory --- Why Passive Pattern Works;
6.~Constrained Control Law; 7.~Theoretical Results;
8.~Estimation of Load / Volatility (Operationalisation);
9.~Experimental Design and Reporting Template;
10.~Ethics, Autonomy, and Safety Guarantees; 11.~Discussion;
12.~Limitations; 13.~Conclusion; Appendices.

\subsection*{Abstract}

We present \emph{Passive Pattern}, a mathematically defined
framework for stability-first dialogue control under non-ergodic
uncertainty. Dialogue is modelled as a partially observable
controlled stochastic process with an absorbing collapse set. Unlike
outcome-optimising interaction models, Passive Pattern optimises
\emph{interactional continuity} by minimising:

\begin{enumerate}
\item cognitive / relational load,
\item volatility (variance of interaction state), and
\item irreversibility (commitment that cannot be safely undone),
\end{enumerate}

\noindent
subject to hard autonomy constraints. This version contributes a
mechanistic theory explaining why the protocol reduces collapse
risk:

\begin{enumerate}[label=(\alph*)]
\item many dialogue breakdowns are variance-amplification phenomena
under partial observability;
\item load increases shrink the feasible action set, increasing
sensitivity to error; and
\item persuasive / commitment-heavy moves increase irreversibility,
which accelerates tail-risk in absorbing systems.
\end{enumerate}

\noindent
We derive a constrained control law, provide stability and safety
results under mild assumptions, and specify an experimental design
and reporting template suitable for peer review.

\subsection{Introduction}
\label{sec:pp-introduction}

Human dialogue is fragile: a single escalation can permanently
terminate cooperation. In such systems, maximising short-run
progress can be dominated by the risk of catastrophic termination.
This is a non-ergodic setting: outcomes are governed by time-path
survival, not ensemble expectations.

The core thesis of this paper is that a large class of real-world
conversations should be treated as \emph{interactive risk-control
problems} rather than persuasion or task-completion problems.
Passive Pattern provides a control-theoretic formulation and a
protocol implementing the resulting control law.

\subsection{Claims and Scope}
\label{sec:pp-claims-scope}

\paragraph{Primary claims (strong form).}

\begin{enumerate}
\item \textbf{(Control formulation)} Dialogue stability can be
formulated as constrained stochastic control with an absorbing
collapse set.
\item \textbf{(Mechanism)} Escalation is often driven by variance
amplification under partial observability; silence / deferral / exit
act as variance-damping operators.
\item \textbf{(Feasibility)} Load increase shrinks the feasible
action set, raising collapse sensitivity; reducing load expands
controllability.
\item \textbf{(Irreversibility)} Commitment-heavy or persuasive
moves increase irreversibility, which accelerates tail risk in
absorbing systems.
\item \textbf{(Guarantees)} Under mild assumptions, Passive Pattern
yields bounded expected volatility and an upper bound on collapse
probability.
\end{enumerate}

\paragraph{Scope.}
Cooperative or potentially cooperative dialogue where continuity
matters (mentorship, negotiation, service recovery, relationships).

\paragraph{Not assumed.}
Shared goals, perfect rationality, or symmetric information.

\subsection{Related Work}
\label{sec:pp-related-work}

\begin{table}[ht]
\centering
\small
\begin{tabularx}{\linewidth}{lXXX}
\toprule
\textbf{Area} & \textbf{Examples} & \textbf{Objective} & \textbf{How Passive Pattern differs} \\
\midrule
Persuasion / Argumentation
& Argumentation systems, persuasion games
& Change beliefs / choices
& PP does not optimise belief change; it optimises stability with
minimal irreversibility. \\

Motivational Interviewing
& Clinical counselling protocols
& Resolve ambivalence, evoke change-talk
& Therapeutic framing; PP is a general control framework allowing
explicit exit / silence as optimal. \\

Task-oriented Dialogue
& Dialogue management, RL for task completion
& Maximise task success
& PP targets collapse avoidance under non-ergodicity; task success
is secondary or optional. \\

Conversation Analysis
& Turn-taking / repair
& Describe structure
& PP prescribes actions from a constrained control law. \\
\bottomrule
\end{tabularx}
\caption[Related work]{Comparison of Passive Pattern with adjacent
fields. PP is unique in formulating dialogue as constrained
stochastic control with an absorbing collapse set.}
\label{tab:pp-related-work}
\end{table}

\subsection{Setup --- Dialogue as Non-Ergodic Controlled Process}
\label{sec:pp-setup}

\subsubsection{State, Actions, and Collapse}
\label{sec:pp-state-actions-collapse}

Let the latent interaction state be $s_t \in \mathcal{S}$ and the
observable features be $o_t \in \mathcal{O}$. We model dialogue as
a partially observable controlled stochastic process:

\begin{equation}
s_{t+1} \sim P(\cdot \mid s_t, a_t), \qquad
o_t \sim \Phi(\cdot \mid s_t).
\label{eq:pp-state-dynamics}
\end{equation}

We define an absorbing collapse set $\mathcal{F} \subset \mathcal{S}$
such that:

\begin{equation}
\forall s \in \mathcal{F},\quad
\Pr(s' \in \mathcal{F} \mid s, a) = 1.
\label{eq:pp-absorbing}
\end{equation}

Collapse captures irreversible breakdown: hard disengagement,
conflict escalation, trust rupture, or operational failure to
continue. Let $\tau_F = \inf\{t : s_t \in \mathcal{F}\}$.

\subsubsection{Sufficient Control State (Operational State)}
\label{sec:pp-control-state}

Passive Pattern uses an operational state vector $x_t = (L_t, V_t,
E_t)$, where:

\begin{itemize}
\item $L_t$ = load (cognitive + emotional + relational, combined);
\item $V_t$ = volatility (variance / instability of interaction
trajectory);
\item $E_t$ = engagement (probability of continued cooperative
interaction).
\end{itemize}

\noindent
These are estimators or summaries of the latent state $s_t$ and may
be computed from $o_t$ and history.

\subsubsection{Irreversibility}
\label{sec:pp-irreversibility}

We define an irreversibility functional $I(a_t, x_t) \ge 0$
capturing the degree to which action $a_t$ commits the interaction
to a narrow future path (e.g., forcing a decision, escalating
stakes, introducing urgency, or requiring a yes/no answer). High
irreversibility reduces the possibility of returning to neutral
interaction states.

\subsection{Mechanistic Theory --- Why Passive Pattern Works}
\label{sec:pp-mechanistic-theory}

This section provides the missing ``why'': the mechanism by which
the protocol reduces collapse risk. We connect three interacting
mechanisms: variance amplification, feasible-action-set shrinkage,
and irreversibility-driven tail risk.

\subsubsection{Variance Amplification under Partial Observability}
\label{sec:pp-variance-amplification}

In many conversations, the system is partially observable: we cannot
directly measure the true latent state (intentions, boundary,
fatigue, hidden constraints). Under partial observability,
aggressive probing or high-pressure clarification can increase
estimation error and state variance.

We model volatility dynamics (in expectation) as:

\begin{equation}
\mathbb{E}[V_{t+1} \mid x_t, a_t] = \rho V_t + \Delta_V(a_t, x_t),
\label{eq:pp-volatility-dynamics}
\end{equation}

with $\rho \in [0, 1)$ capturing natural dissipation and $\Delta_V$
capturing variance injected by action. Actions like ``force
decision'' or repeated probing often yield $\Delta_V > 0$; stabilising
actions yield $\Delta_V \le 0$.

\paragraph{Variance-damping operators.}
Define the subset

\begin{equation}
\mathcal{A}^{\downarrow} = \{a : \Delta_V(a, x) \le 0
\text{ for relevant } x\}.
\label{eq:pp-variance-damping}
\end{equation}

Passive Pattern explicitly includes silence, deferral, and polite
exit in $\mathcal{A}^{\downarrow}$. These operators reduce variance
by removing injected pressure and allowing the latent state to
relax.

\subsubsection{Feasible-Action-Set Shrinkage (Why Load Matters)}
\label{sec:pp-feasible-action-set}

Load is not merely ``how someone feels''; it is a control-theoretic
proxy for controllability. As load rises, the set of actions that
keep the system within safe bounds shrinks. Let
$\mathcal{A}_{\mathrm{safe}}(x)$ be the set of actions satisfying
safety constraints at state $x$:

\begin{equation}
\mathcal{A}_{\mathrm{safe}}(x) = \{a \in \mathcal{A} : \Pr(\tau_F \le T \mid x, a) \le \varepsilon\}.
\label{eq:pp-safe-action-set}
\end{equation}

We assume a monotonicity condition (supported by operational
observation): higher load reduces safety margins:

\begin{equation}
L' \ge L \Rightarrow
\mathcal{A}_{\mathrm{safe}}(L', V, E) \subseteq
\mathcal{A}_{\mathrm{safe}}(L, V, E).
\label{eq:pp-monotonicity}
\end{equation}

Thus, interventions that reduce $L$ increase the feasible safe
control set, improving controllability and reducing collapse risk.

\subsubsection{Irreversibility as Tail-Risk Accelerator}
\label{sec:pp-irreversibility-tail-risk}

In absorbing systems, the long-run objective is dominated by tail
events. Irreversible actions reduce recovery pathways and amplify
tail risk. Let $p_{\mathrm{fail}}(x, a)$ denote one-step contribution
to collapse risk. We posit:

\begin{equation}
\frac{\partial p_{\mathrm{fail}}(x, a)}{\partial I(a, x)} \ge 0,
\label{eq:pp-tail-risk}
\end{equation}

\noindent
i.e., higher irreversibility increases collapse probability,
especially when $E$ is low or $V$ is high.

This yields a structural critique of persuasion / urgency tactics
in non-ergodic interaction: even if they increase short-run
``progress'' (e.g., forced decision), they increase irreversibility
and therefore tail risk. Passive Pattern reduces irreversibility by
preferring low-commitment actions and allowing safe withdrawal.

\subsubsection{Why Silence / Deferral / Exit Improve Long-Horizon
Outcomes}
\label{sec:pp-silence-improves}

Combining the above:

\begin{itemize}
\item Variance damping reduces $V_t$ (lower escalation probability).
\item Load reduction expands $\mathcal{A}_{\mathrm{safe}}$ (more
controllable region).
\item Irreversibility minimisation reduces tail-risk acceleration.
\end{itemize}

\noindent
Even if these moves reduce immediate ``progress,'' they improve the
probability of continued interaction, enabling future progress
without crossing collapse boundaries. In non-ergodic systems, this
is rational: preserving the process dominates maximising any single
step.

\subsection{Constrained Control Law}
\label{sec:pp-constrained-control-law}

\subsubsection{Objective}
\label{sec:pp-objective}

We minimise a risk-sensitive cost with explicit irreversibility
penalty:

\begin{equation}
J(\pi) = \mathbb{E}_\pi\!\left[\sum_{t=0}^T
\bigl(\alpha L_t + \beta V_t + \gamma I(a_t, x_t)\bigr)\right].
\label{eq:pp-objective}
\end{equation}

\subsubsection{Safety and Autonomy Constraints}
\label{sec:pp-safety-autonomy}

\paragraph{Safety (chance constraint).}

\begin{equation}
\Pr_\pi(\tau_F \le T) \le \varepsilon.
\label{eq:pp-chance-constraint}
\end{equation}

\paragraph{Autonomy.}
The autonomy constraint set $\mathcal{A}_{\mathrm{aut}}(x)$ excludes
force, deception, urgency, and shaping acts:

\begin{equation}
a_t \in \mathcal{A}_{\mathrm{aut}}(x_t) \subset \mathcal{A}.
\label{eq:pp-autonomy}
\end{equation}

\subsubsection{Control Policy (Canonical Implementation)}
\label{sec:pp-control-policy}

\begin{lstlisting}[basicstyle=\ttfamily\small,frame=single]
Inputs:
  estimates (L_t, V_t, E_t)
  risk model p_fail(x, a)
  thresholds (epsilon, L_hi, V_hi)

Step 1: compute risk for safe actions:
  R(a) = P(tau_F <= T | x_t, a)

Step 2:
  if min_a R(a) > epsilon:
    choose EXIT (polite) or SILENCE
      (if exit is socially harmful)
  elif L_t >= L_hi or V_t >= V_hi:
    choose a in A_down (variance-damping):
      SILENCE or DEFER or REFLECT
  else:
    choose low-commitment INQUIRY or SUMMARY
      minimizing alpha*L + beta*V + gamma*I

Step 3:
  update estimators
  terminate calibration when stability condition holds
\end{lstlisting}

\subsubsection{Variable-Length Calibration (Shortest Route)}
\label{sec:pp-variable-length}

Calibration continues only until the minimum required uncertainty
reduction is achieved. Define a confidence score $C_t$ on the
estimated safe set and a stability score $S_t$. Stop calibration at
time $t$ if $C_t \ge \overline{C}$ and $S_t \ge \overline{S}$.

\subsection{Theoretical Results}
\label{sec:pp-theoretical-results}

We provide results under standard stochastic control assumptions:
bounded costs, existence of safe actions in a neighbourhood of
stable states, and correct specification of autonomy constraints.
Proofs are provided as sketches in Appendix A and can be expanded
into full measure-theoretic detail if required by venue.

\subsubsection{Assumptions}
\label{sec:pp-assumptions}

\begin{description}
\item[(A1)] \textbf{Absorbing collapse}: $\mathcal{F}$ is absorbing.
\item[(A2)] \textbf{Variance damping actions exist}: For relevant
$x$, $\mathcal{A}^{\downarrow}(x) \cap
\mathcal{A}_{\mathrm{aut}}(x) \ne \emptyset$.
\item[(A3)] \textbf{Monotone feasibility}: $L$ increases shrink
$\mathcal{A}_{\mathrm{safe}}$ as in
Section~\ref{sec:pp-feasible-action-set}.
\item[(A4)] \textbf{Risk estimator conservative}: Estimated
$\widehat{p}_{\mathrm{fail}}(x, a)$ upper-bounds true
$p_{\mathrm{fail}}(x, a)$ with high probability.
\end{description}

\subsubsection{Lemma 1 (Variance Contraction under Damping)}
\label{sec:pp-lemma-1}

\begin{lemma}[Variance Contraction under Damping]
\label{lem:pp-variance-contraction}
If $a_t \in \mathcal{A}^{\downarrow}$, then for relevant $x_t$,

\begin{equation}
\mathbb{E}[V_{t+1} \mid x_t, a_t] \le \rho V_t,
\quad \text{for some } \rho \in [0, 1),
\end{equation}

\noindent
i.e., expected volatility contracts.
\end{lemma}

\subsubsection{Lemma 2 (Feasible-Set Expansion under Load
Reduction)}
\label{sec:pp-lemma-2}

\begin{lemma}[Feasible-Set Expansion under Load Reduction]
\label{lem:pp-feasible-expansion}
If an action sequence reduces load from $L$ to $L' < L$ while not
increasing $V$, then:

\begin{equation}
\mathcal{A}_{\mathrm{safe}}(L, V, E) \subseteq
\mathcal{A}_{\mathrm{safe}}(L', V, E),
\end{equation}

\noindent
expanding the safe controllability region.
\end{lemma}

\subsubsection{Theorem 1 (Safety: Chance Constraint Satisfaction)}
\label{sec:pp-thm-1}

\begin{theorem}[Safety: Chance Constraint Satisfaction]
\label{thm:pp-safety}
Under (A1)--(A4), the Passive Pattern policy that selects actions
only from $\mathcal{A}_{\mathrm{safe}} \cap
\mathcal{A}_{\mathrm{aut}}$ satisfies:

\begin{equation}
\Pr_\pi(\tau_F \le T) \le \varepsilon.
\end{equation}
\end{theorem}

\subsubsection{Theorem 2 (Stability: Bounded Expected Volatility)}
\label{sec:pp-thm-2}

\begin{theorem}[Stability: Bounded Expected Volatility]
\label{thm:pp-stability}
Under (A1)--(A3) and bounded injected variance outside
$\mathcal{A}^{\downarrow}$, the policy yields:

\begin{equation}
\sup_{t \le T} \mathbb{E}[V_t] \le V_0 +
\frac{\sup_{x, a} \Delta_V^+(a, x)}{1 - \rho},
\end{equation}

\noindent
where $\Delta_V^+ = \max(\Delta_V, 0)$.
\end{theorem}

\subsection{Estimation of Load / Volatility (Operationalisation)}
\label{sec:pp-estimation}

Passive Pattern requires estimators $\widehat{L}_t$,
$\widehat{V}_t$, $\widehat{E}_t$ computed from observable signals
and dialogue history. This section specifies concrete estimators;
venues may treat them as engineering choices.

\subsubsection{Load Estimator}
\label{sec:pp-load-estimator}

Define a feature vector $z_t$ (speech rate, hesitation markers,
response latency, sentiment variance, boundary-language). A simple
estimator is:

\begin{equation}
\widehat{L}_t = w_L^\top z_t,
\label{eq:pp-load-estimator}
\end{equation}

\noindent
with weights $w_L$ calibrated on labelled data or expert rules.

\subsubsection{Volatility Estimator}
\label{sec:pp-volatility-estimator}

Let $y_t$ be an interaction quality score (e.g., cooperative tone,
agreement markers). Define volatility as an exponential moving
variance:

\begin{equation}
\widehat{V}_t = (1 - \lambda) \widehat{V}_{t-1} +
\lambda (y_t - \overline{y}_t)^2.
\label{eq:pp-volatility-estimator}
\end{equation}

\subsubsection{Engagement Estimator}
\label{sec:pp-engagement-estimator}

Engagement is modelled as a hazard process:

\begin{equation}
\widehat{E}_t = \sigma(w_E^\top z_t),
\label{eq:pp-engagement-estimator}
\end{equation}

\noindent
where $\sigma$ is the logistic function.

\subsection{Experimental Design and Reporting Template}
\label{sec:pp-experimental-design}

To satisfy peer review, experiments must (i) define baselines, (ii)
report collapse metrics, and (iii) show robustness across contexts.
Because empirical datasets are venue-dependent, we provide a
reporting template and recommended minimal experiments.

\subsubsection{Baselines}
\label{sec:pp-baselines}

\begin{description}
\item[Baseline A (Direct)] Direct questioning / direct proposal
loops.
\item[Baseline B (Clarification-heavy)] Repeated clarification and
commitment checks.
\item[Baseline C (Task-optimised)] Maximise task completion (if
available).
\end{description}

\subsubsection{Metrics}
\label{sec:pp-metrics}

\begin{itemize}
\item \textbf{Breakdown rate}: $\Pr(\tau_F \le T)$.
\item \textbf{Time-to-breakdown}: $\mathbb{E}[\tau_F]$ (censored).
\item \textbf{Volatility integral}: $\sum_t \widehat{V}_t$.
\item \textbf{Irreversibility exposure}: $\sum_t I(a_t, x_t)$.
\item \textbf{Autonomy violations}: count of excluded actions
attempted (should be zero).
\end{itemize}

\subsubsection{Minimal Study (Human-in-the-loop)}
\label{sec:pp-minimal-study}

\begin{enumerate}
\item Randomly assign interaction scripts to Passive Pattern vs
baseline.
\item Use identical goal prompts; allow exit at any time.
\item Collect load ratings, disengagement, and conflict escalation
markers.
\end{enumerate}

\paragraph{Note on numbers.}
This document avoids fabricating empirical results. Tables should
be populated with actual measurements from the target deployment or
study.

\subsection{Ethics, Autonomy, and Safety Guarantees}
\label{sec:pp-ethics-autonomy}

Passive Pattern enforces non-manipulation via hard constraints on
the action set. This is stronger than a guideline: prohibited
actions are not available to the optimiser.

\paragraph{Autonomy constraint.}

\begin{equation}
\mathcal{A}_{\mathrm{aut}}(x) = \mathcal{A} \setminus
\{\mathrm{force}, \mathrm{deceive}, \mathrm{urgency},
\mathrm{shape}\}.
\label{eq:pp-autonomy-set}
\end{equation}

\paragraph{Withdrawal as safe outcome.}
Withdrawal is a legitimate terminal action and is treated as a safe,
successful outcome for system continuity.

\subsection{Discussion}
\label{sec:pp-discussion}

The mechanistic theory explains why Passive Pattern generalises: it
is not a set of clever questions, but a control framework that
dampens variance, expands controllability, and avoids
irreversibility-driven tail risk. This makes it compatible with
stability-first decision systems (e.g., NRMO+Z) and with
conservative operational environments.

\subsection{Limitations}
\label{sec:pp-limitations}

\begin{itemize}
\item Estimators may be misspecified; conservative risk bounds are
recommended (A4).
\item In adversarial settings, exit / silence may be exploited;
combine with NRMO governance.
\item Task completion is not guaranteed; Passive Pattern prioritises
continuity over completion.
\end{itemize}

\subsection{Conclusion}
\label{sec:pp-conclusion}

We provided a full mechanistic theory and a constrained stochastic
control formulation for stability-first dialogue. Passive Pattern
reduces collapse risk not by persuasion but by variance damping,
feasible-action-set expansion, and irreversibility avoidance under
non-ergodicity.

\subsection*{Appendix A --- Proof Sketches}
\label{sec:pp-app-A}

\paragraph{Lemma 1 (Variance Contraction)} follows from defining
$\mathcal{A}^{\downarrow}$ as the set of actions with non-positive
variance injection $\Delta_V$ and using contraction $\rho < 1$.

\paragraph{Theorem 1 (Safety)} follows by selecting only
conservative-safe actions under an upper-bounding risk estimator
(A4), ensuring the chance constraint holds.

\paragraph{Theorem 2 (Stability)} follows from bounding the positive
part of variance injection and summing the resulting geometric
series with contraction factor $\rho$.

\paragraph{Lemma 2 (Feasible-Set Expansion)} If an action sequence
reduces load from $L$ to $L' < L$ while not increasing $V$, then:

\begin{equation}
\mathcal{A}_{\mathrm{safe}}(L, V, E) \subseteq
\mathcal{A}_{\mathrm{safe}}(L', V, E),
\end{equation}

\noindent
expanding the safe controllability region.

\subsection*{Appendix B --- Reference Implementation (Pseudocode)}
\label{sec:pp-app-B}

\begin{lstlisting}[basicstyle=\ttfamily\scriptsize,frame=single,language=Python]
class PassivePatternController:
    def __init__(self, epsilon, L_hi, V_hi,
                 phi, risk_model, autonomy_filter):
        self.epsilon         = epsilon
        self.L_hi            = L_hi
        self.V_hi            = V_hi
        self.phi             = phi
        self.risk_model      = risk_model
        self.autonomy_filter = autonomy_filter

    def choose_action(self, x, candidate_actions):
        A = [a for a in candidate_actions
             if self.autonomy_filter(a, x)]
        risks = {a: self.risk_model(x, a) for a in A}

        if min(risks.values()) > self.epsilon:
            return "EXIT"
            # or "SILENCE" if social context penalises exit

        if x.L >= self.L_hi or x.V >= self.V_hi:
            return self._variance_damping(A, x)

        return self._low_commitment(A, x)

    def _variance_damping(self, A, x):
        for a in ["SILENCE", "DEFER", "REFLECT"]:
            if a in A:
                return a
        return "SUMMARY"

    def _low_commitment(self, A, x):
        for a in ["INQUIRY", "REFLECT", "SUMMARY"]:
            if a in A:
                return a
        return "SILENCE"
\end{lstlisting}

\subsection*{Appendix C --- Integration with NRMO+Z}
\label{sec:pp-app-C}

Passive Pattern is an interpersonal control layer compatible with
NRMO+Z: NRMO constrains action selection under non-ergodicity,
while Passive Pattern constrains dialogue actions under non-ergodic
relational risk. Both share the principles of:

\begin{enumerate}
\item absorbing failure avoidance,
\item chance constraints, and
\item reversibility preference.
\end{enumerate}

\paragraph{Version.}
Passive Pattern v1.5 (Mechanistic Theory) --- operational.

\subsection*{Appendix D --- Integration Constraint with NRMO /
Parallel OODA}
\label{sec:pp-app-D}

When integrated into NRMO-based decision systems, Passive Pattern
must be treated as a \emph{protective control layer}, not as a
decision-making or optimisation module. Specifically:

\begin{itemize}
\item Passive Pattern may signal risk, suggest stabilisation, or
request deferral, suspension, or exit.
\item It must \textbf{not} generate preferences, rankings,
strategies, or action plans.
\item Outputs of Passive Pattern are limited to \emph{triggers and
constraints}, not decisions.
\end{itemize}

\paragraph{Interfaces.}
In NRMO architectures, Passive Pattern primarily interfaces with:

\begin{itemize}
\item Passive Pattern $\to$ Parallel OODA (Observe / Orient).
\item Passive Pattern $\to$ Passive Pattern execution
(silence, deferral, exit triggers).
\end{itemize}

\paragraph{Defensive constraint-based intervention.}
In exceptional defensive conditions, Passive Pattern may impose
constraint-based interventions on Decide or Act stages, such as
halting, delaying, or narrowing available actions. \textbf{Passive
Pattern must never lead, optimise, or persuasively shape decisions
or actions.} Final decision authority always remains with NRMO.
